% Part: applied-modal-logics
% Chapter: temporal-logic
% Section: possible-histories

\documentclass[../../../include/open-logic-section]{subfiles}

\begin{document}

\olfileid{aml}{tl}{poss}

\olsection{في التواريخ الممكنة}

لما كانت النماذج العلائقية التي تقدمت للمنطق الزمني لا تقيد علاقة النفاذ،
اتسعت لضروب كثيرة من البنية، ومنها تفرع الماضي والمستقبل. غير أنا قد نريد
تمثيل هذا التفرع تمثيلًا أصرح بالجهة، فنجعل متتاليات الحوادث تواريخ ممكنة.
ولا يوجب ذلك تغيير اللغة نفسها؛ وإن جاز أن نلحق بها المؤثرين الجهويين
المعهودين~$\Box$ و$\Diamond$، أو نزيد علاقات نفاذ معرفية تمثل تغير معرفة
الفاعلين على مر الزمن.

\begin{defn}
  \emph{نموذج التواريخ الممكنة} للغة الزمنية هو الثلاثي
  $\mModel{M}=\tuple{\OLId{latin-looped}{T},\OLId{latin-looped}{C},\OLId{latin-looped}{V}}$، وفيه:
  \begin{enumerate}
    \item $\OLId{latin-looped}{T}$ مجموعة غير فارغة، تُفسر عناصرها بوصفها حالات في الزمن.
    \item $\OLId{latin-looped}{C}$ مجموعة من مسارات التنفيذ، أي \emph{التواريخ الممكنة}
      لنظام ما؛ وبعبارة أخرى،~$\OLId{latin-looped}{C}$ مجموعة من المتتاليات~$\sigma$ للحالات
      $\OLId{latin-ordinary}{s}_\OLMathNumeral{1}$، $\OLId{latin-ordinary}{s}_\OLMathNumeral{2}$،~$\OLId{latin-ordinary}{s}_\OLMathNumeral{3}$، \dots، حيث كل~$\OLId{latin-ordinary}{s}_\OLId{latin-ordinary}{i} \in \OLId{latin-looped}{T}$.
    \item $\OLId{latin-looped}{V}$ دالة تسند إلى كل~!!{propositional
      variable}~$\OLId{latin-ordinary}{p}$ مجموعة~$\OLId{latin-looped}{V}(\OLId{latin-ordinary}{p})$ من النقاط الزمنية.
  \end{enumerate}
  ونفترض، طلبًا للتيسير، أن كل تاريخ في~$\OLId{latin-looped}{C}$ تكون جميع لواحقه فيها أيضًا.
  فإذا كانت~$\OLId{latin-ordinary}{s}_\OLMathNumeral{1}$، $\OLId{latin-ordinary}{s}_\OLMathNumeral{2}$،~$\OLId{latin-ordinary}{s}_\OLMathNumeral{3}$ متتالية في~$\OLId{latin-looped}{C}$، كانت~$\OLId{latin-ordinary}{s}_\OLMathNumeral{2}$، $\OLId{latin-ordinary}{s}_\OLMathNumeral{3}$
  وكذلك~$\OLId{latin-ordinary}{s}_\OLMathNumeral{3}$ متتاليتين فيها. وإذا اجتمعت حالتان~$\OLId{latin-ordinary}{s}_\OLId{latin-ordinary}{i}$ و$\OLId{latin-ordinary}{s}_\OLId{latin-ordinary}{j}$ في
  متتالية~$\sigma$، قلنا إن~$\OLId{latin-ordinary}{s}_\OLId{latin-ordinary}{i} \prec_\sigma \OLId{latin-ordinary}{s}_\OLId{latin-ordinary}{j}$ متى كان $\OLId{latin-ordinary}{i}<\OLId{latin-ordinary}{j}$.
  وإذا كان~$\OLId{latin-ordinary}{t} \in \OLId{latin-looped}{V}(\OLId{latin-ordinary}{p})$ قلنا إن~$\OLId{latin-ordinary}{p}$ \emph{صادق عند}~$\OLId{latin-ordinary}{t}$.
\end{defn}

وإنما يتغير التقويم من جهة واحدة: فصدق !!a{formula} عند النقطة الزمنية~$\OLId{latin-ordinary}{t}$
في النموذج~$\mModel M$ يقوّم الآن بالإضافة إلى تاريخ~$\sigma$ تظهر فيه~$\OLId{latin-ordinary}{t}$
حالةً. وأما دلالات ال!!{propositional variable}s والروابط الدالية الصدق
فباقية على ما في~\olref[sem]{defn:tmodels}، إذ لا يشير شيء منها إلى
$\sigma$. والذي نعيد حده هو مؤثر المستقبل~$\Ftemp$، ثم نضيف المؤثر
$\Diamond$ بالنظر إلى هذه التواريخ.

\begin{defn}\ollabel{defn:phmodels}
  \emph{صدق !!a{formula}~$!A$ عند~$\OLId{latin-ordinary}{t},\sigma$} في~$\mModel M$، ورمزه
  $\mSat{\OLId{latin-looped}{M}}{!A}[\OLId{latin-ordinary}{t},\sigma]$:
  \begin{enumerate}
  \item\ollabel{defn:sub:phmodels-f} \indcase{!A}{\Ftemp
    !B}{$\mSat{\OLId{latin-looped}{M}}{\indfrm}[\OLId{latin-ordinary}{t},\sigma]$ إذا وفقط إذا كان
    $\mSat{\OLId{latin-looped}{M}}{!B}[\OLId{latin-ordinary}{t}',\sigma]$ لبعض~$\OLId{latin-ordinary}{t}' \in \OLId{latin-looped}{T}$ بحيث
    $\OLId{latin-ordinary}{t} \prec_\sigma \OLId{latin-ordinary}{t}'$.}
  \item\ollabel{defn:sub:phmodels-diamond} \indcase{!A}{\Diamond
    !B}{$\mSat{\OLId{latin-looped}{M}}{\indfrm}[\OLId{latin-ordinary}{t},\sigma]$ إذا وفقط إذا كان
    $\mSat{\OLId{latin-looped}{M}}{!B}[\OLId{latin-ordinary}{t},\sigma']$ لبعض~$\sigma' \in \OLId{latin-looped}{C}$ تظهر فيه~$\OLId{latin-ordinary}{t}$.}
  \end{enumerate}
\end{defn}

وعلى هذا القياس تحد مؤثرات أخرى زمنية وجهوية. وقد صار في وسعنا أن نمثل
قضية تجمع الزمان والجهة؛ فقولنا «لن يقع~$\OLId{latin-ordinary}{p}$ في هذا التاريخ، مع إمكان وقوعه
في تاريخ بديل» نرمز إليه بالصيغة~$\lnot \Ftemp \OLId{latin-ordinary}{p} \land \Diamond \Ftemp \OLId{latin-ordinary}{p}$. وهي تصدق عند
نقطة في تاريخ لا يصير فيه~$\OLId{latin-ordinary}{p}$ صادقًا في حالة آتية، إذا وجد تاريخ بديل
يصير فيه~$\OLId{latin-ordinary}{p}$ صادقًا في المستقبل.

\end{document}
